\documentclass[11pt,a4paper]{article}

\usepackage[utf8]{inputenc}
\usepackage[indonesian]{babel}
\usepackage{amsmath, amssymb, amsfonts}
\usepackage{geometry}
\geometry{margin=1in, headheight=14pt}
\usepackage{graphicx}
\usepackage{booktabs}
\usepackage{tcolorbox}
\tcbuselibrary{skins,breakable}
\usepackage{enumitem}
\usepackage{fancyhdr}

\definecolor{unibblue}{RGB}{0, 51, 102}
\definecolor{unibgold}{RGB}{212, 175, 55}
\definecolor{darkgreen}{RGB}{34, 139, 34}

\pagestyle{fancy}
\fancyhf{}
\rhead{\textbf{TKE-41XX: Optimisasi untuk Teknik Elektro}}
\lhead{Lembar Kerja Mahasiswa (LKM) Minggu 2}
\rfoot{Halaman \thepage}

\title{\vspace{-1cm}\textbf{LEMBAR KERJA MAHASISWA (LKM) MINGGU KE-2:}\\Optimisasi Linier (LP) dan Metode Grafis 2D}
\author{Program Studi S1 Teknik Elektro --- Universitas Bengkulu}
\date{Semester Ganjil 2026}

\begin{document}

\maketitle

\begin{tcolorbox}[colback=unibblue!5,colframe=unibblue,title=\textbf{Identitas Mahasiswa dan Instruksi}]
\begin{tabular}{ll}
\textbf{Nama Mahasiswa} : & \hrulefill\hspace{6cm} \\
\textbf{NPM} : & \hrulefill\hspace{6cm} \\
\textbf{Kelompok / Kelas} : & \hrulefill\hspace{6cm} \\
\end{tabular}
\vspace{0.2cm}
\begin{itemize}[leftmargin=*]
    \item Kerjakan LKM ini secara kolaboratif dalam kelompok (2--3 mahasiswa).
    \item Jawablah setiap pertanyaan secara sistematis pada ruang yang disediakan.
    \item Soal ini mencakup ranah kognitif Taksonomi Bloom tingkat **C1 hingga C6**.
\end{itemize}
\end{tcolorbox}

\vspace{0.3cm}

\section*{Bagian 1: Asumsi Dasar LP dan Geometri Poligon (C1--C2)}

\begin{enumerate}[label=\textbf{Soal \arabic*:}, leftmargin=*]
    \item \textbf{[C1 --- Mengingat]} Jelaskan 4 asumsi mendasar dalam Linear Programming (Proposionalitas, Aditivitas, Divisibilitas, Kepastian)!
    
    \begin{tcolorbox}[colback=white,colframe=gray!50,height=3.5cm,title=Ruang Jawaban Soal 1]
    \end{tcolorbox}

    \item \textbf{[C2 --- Memahami]} Mengapa solusi optimal pada masalah LP terkendala poligon dijamin selalu berada pada titik pojok ekstrim (*Corner Point Feasible*)?
    
    \begin{tcolorbox}[colback=white,colframe=gray!50,height=3.5cm,title=Ruang Jawaban Soal 2]
    \end{tcolorbox}
\end{enumerate}

\section*{Bagian 2: Solusi Grafis dan Pemrograman Python (C3--C4)}

\begin{enumerate}[label=\textbf{Soal \arabic*:}, leftmargin=*, resume]
    \item \textbf{[C3 --- Mengaplikasikan]} Diberikan model LP:
    \[
        \max Z = 3 x_1 + 5 x_2 \quad \text{s.t.} \quad x_1 + 2 x_2 \le 12, \quad 3 x_1 + 2 x_2 \le 18, \quad x_1, x_2 \ge 0
    \]
    Gambarkan daerah feasibel secara grafis dan tentukan nilai $Z$ pada seluruh titik pojok ekstrim!
    
    \begin{tcolorbox}[colback=white,colframe=gray!50,height=5cm,title=Ruang Jawaban Soal 3]
    \end{tcolorbox}

    \item \textbf{[C4 --- Menganalisis]} Tuliskan skrip Python menggunakan `scipy.optimize.linprog` untuk memverifikasi hasil titik pojok dari Soal 3!
    
    \begin{tcolorbox}[colback=white,colframe=gray!50,height=5.5cm,title=Ruang Jawaban Soal 4]
    \end{tcolorbox}
\end{enumerate}

\section*{Bagian 3: Evaluasi Feasibilitas dan Perancangan Kasus (C5--C6)}

\begin{enumerate}[label=\textbf{Soal \arabic*:}, leftmargin=*, resume]
    \item \textbf{[C5 --- Mengevaluasi]} Jika kendala ketiga $x_1 + x_2 \ge 15$ ditambahkan pada Soal 3, evaluasi bagaimana perubahan ruang solusi poligon dan status feasibilitas masalah tersebut!
    
    \begin{tcolorbox}[colback=white,colframe=gray!50,height=4.5cm,title=Ruang Jawaban Soal 5]
    \end{tcolorbox}

    \item \textbf{[C6 --- Mencipta]} Buatlah sebuah masalah *Linear Programming* 2D riil dalam sistem jaringan distribusi listrik (alokasi beban trafo / pembangkit terdistribusi) lengkap dengan garis isoprofit dan titik pojok optimalnya!
    
    \begin{tcolorbox}[colback=white,colframe=gray!50,height=6.5cm,title=Ruang Jawaban Soal 6]
    \end{tcolorbox}
\end{enumerate}

\end{document}
